\section{Введение}

Перевернутый маятник является классической задачей теории управления, демонстрирующей возможность стабилизации неустойчивых систем с помощью обратной связи. Несмотря на простоту механической конструкции, эта система обладает сложной динамикой и служит эталоном для тестирования методов управления.

\subsection{Цель работы}
Исследование возможности стабилизации перевернутого маятника в верхнем положении с помощью горизонтального перемещения точки опоры. Анализ различных стратегий управления.

\subsection{Постановка задачи}
Рассматривается математический маятник массой \(m\) и длиной \(l\), точка подвеса которого может перемещаться горизонтально. Управлением служит ускорение точки опоры \(\ddot{u}\). Необходимо найти закон управления, обеспечивающий асимптотическую устойчивость верхнего положения равновесия (\(\varphi = 0, \dot{\varphi} = 0, u = 0, \dot{u} = 0\)).